\documentclass[a4paper]{article} %

\date{} %

% Please, do not change any of the above lines

\usepackage{amssymb} % add other LaTeX packages you need here.

\begin{document}
\setcounter{page}{1} % Please, do not delete this line

% Your title
\title{Universal algebraic geometry of normal bands}

% Authors
\author{Kristóf Varga, Tamás Waldhauser \\ %
       University of Szeged\\ \\ % Affiliation 1
       % Author 4 \\ % If any other author with different Affilation
       % University of ZZWW\\ \\ % Affiliation 2 (if needed)
       % New author \\
       % New affiliation \\
       % Add authors and affiliation as needed
       \texttt{vkrist99@gmail.com} % Only one corresponding e-mail, please
       }%

\maketitle

% The abstract

Given an algebra $\mathbb A = (A; F )$, we say that a set of $n$-tuples $S \subseteq A^n$ is an ($n$-ary) \emph{algebraic set} if $S$ is the set of solutions of a system of equations over A. 
The collection of all algebraic sets is called the \emph{algebraic geometry} of $\mathbb A$ \cite{P}.
The algebraic geometry of an algebra only depends on the clone of its term operations.
If two clones defined on the same underlying set have the same algebraic geometry, then we say that the two clones are \emph{algebraically equivalent}.
For a fixed algebra $\mathbb A$, two natural problems arise:
\begin{itemize}
    \item characterize the algebraic sets over $\mathbb A$;
    \item characterize the clones on the set $A$ that are algebraically equivalent to the clone of term operations of $\mathbb A$.
\end{itemize}
% Some results of our ongoing research on these two problems will be
% presented in this talk, in the case when $\mathbb A$ is a normal band.

The set of $n$-ary algebraic sets of $\mathbb A$ forms a lattice with respect to inclusion; furthermore, the structure of this lattice is completely determined by the quasivariety generated by $\mathbb A$ if it is a normal band.
Using this connection, we managed to characterize the lattice of $n$-ary algebraic sets of every normal band up to isomorphism.

We obtained a similarly complete answer for the second problem as well: except for some special cases, a clone is algebraically equivalent to the clone of a normal band $\mathbb A$ if and only if it is the clone of $\mathbb A$.
\begin{thebibliography}{99}
\bibitem{P} 
B. Plotkin. 
\newblock \emph{Some results and problems related to universal algebraic geometry,} 
\newblock Internat. J. Algebra Comput. \textbf{17}, no. 5-6, 1133--1164 (2007)
% \bibitem{mundici1986}
% D.~Mundici.
% \newblock \emph{Interpretation of {A}{F} ${C}^*$-algebras in {L}ukasziewicz sentential calculus,}
% \newblock J. Functional Analysis \textbf{65} 15--53 (1986)
% \bibitem{gratzer1998}
% G.~Gr{\" a}tzer.
% \newblock \emph{General Lattice Theory,}
% \newblock Birkh{\" a}user, second edition, 1998.

\end{thebibliography}

\end{document}